\documentclass[12pt]{article}
\usepackage[utf8]{inputenc}
\usepackage[T1]{fontenc}
\usepackage[margin=1in]{geometry}
\usepackage{setspace}
\usepackage{xcolor}
\usepackage[colorlinks=true, urlcolor=blue!60!black, linkcolor=blue!60!black]{hyperref}
\setlength{\parindent}{0em}
\setlength{\parskip}{0.6em}
\renewcommand{\baselinestretch}{1.0}
\begin{document}
\begin{center}
{\LARGE \textit{303 Creative LLC v.\ Elenis}, 600 U.S.\ 570 (2023)\par}
\vspace{0.5em}
{\large Case Brief --- Two-Page Edition\par}
\end{center}

\section{Memory Jogger}
Colorado may not use its public accommodations law to compel a website designer to create custom wedding websites conveying messages she rejects (\href{https://supreme.justia.com/cases/federal/us/600/21-476/}{Justia}).

\section{Detailed Case Facts}
Lorie Smith of 303 Creative LLC wanted to design wedding websites only for opposite-sex weddings, because of her Christian beliefs, and to post a statement saying so. Colorado's Anti-Discrimination Act (CADA) bars public accommodations from refusing service based on sexual orientation or announcing such refusals. The parties stipulated that her websites are ``original, customized'' and ``expressive'' and that she serves LGBTQ+ clients in other projects.

\section{Procedural History}
In Smith's pre-enforcement suit, the district court ruled for Colorado; the Tenth Circuit affirmed 2--1, holding that CADA survived strict scrutiny because her services were ``unique.''

\section{Judicial Votes}
Reversed, 6--3. Gorsuch wrote for Roberts, Thomas, Alito, Kavanaugh, and Barrett; Sotomayor dissented, joined by Kagan and Jackson.

\section{Holding}
A state may not apply its public accommodations law to compel a business that creates customized, expressive ``pure speech'' to convey messages it disagrees with, even if it sells to the public for profit. The Court set no general test for other services.

\section{Analysis of Opinions}
\textbf{Majority:} Under \textit{Barnette}, \textit{Hurley}, and \textit{Dale}, government may not compel a speaker's message. Smith's websites are her own pure speech; profit and LLC status do not matter; her objection is to a message, not a status; and treating unique talent as grounds for coercion would invert the First Amendment. \textbf{Dissent:} CADA regulates the conduct of refusing service based on status, with only an ``incidental'' burden on speech under \textit{O'Brien} and \textit{FAIR}; the Court had long denied businesses a right to discriminate, and redefining the service collapses status into message.

\section{Examples: Future Applications}
\textbf{Same-Side:} a wedding photographer who creates custom visual narratives may decline a same-sex wedding. \textbf{Opposite-Side:} a reception venue renting space and serving food must serve all couples. \textbf{Fence-Sitter:} a custom cake baker or florist, whose partly functional craft may or may not be ``expressive activity.''

\section{Critique}
Scholars fault the majority for never naming a tier of scrutiny or an attribution test, and Robert Post calls its approach an ``impossible abstraction.'' Others argue the stipulated, pre-enforcement record let the opinions ``talk past each other'' and carved a hole in civil rights law.

\section{Key Quotations}
\begin{quote}
``When a state public accommodations law and the Constitution collide, there can be no question which must prevail.''
\end{quote}
\noindent --- Justice Gorsuch, for the Court.

\end{document}
